\chapter{KESIMPULAN DAN SARAN}

\section{Kesimpulan}
Berdasarkan hasil perancangan, implementasi, dan pengujian yang telah dilakukan, dapat disimpulkan beberapa hal sebagai berikut.

\begin{enumerate}
    \item Model CNN berhasil diimplementasikan pada board FPGA menggunakan representasi fixed-point 16-bit Q2.14, yaitu dua bit untuk bagian integer termasuk bit tanda dan 14 bit untuk bagian fraksi. Format ini digunakan karena parameter model tidak keluar dari rentang $-2$ sampai $2$, sehingga resolusi pecahan dapat diprioritaskan. Proses kuantisasi dilakukan terhadap parameter bobot dan bias model sehingga dapat digunakan pada modul Verilog tanpa menggunakan operasi floating-point.

    \item Model Python mencapai akurasi 98,2\% pada data latih dan 97,2\% pada data uji. Pada evaluasi Verilog menggunakan 100 sampel untuk setiap label, kelas prediksi yang dihasilkan sama dengan keluaran model Python untuk sampel masukan yang sama. Perbedaan nilai numerik akibat kuantisasi, pembulatan, pemotongan bit, dan akumulasi operasi MAC tidak mengubah hasil klasifikasi akhir berdasarkan argmax pada sampel yang dievaluasi.

    \item Hasil implementasi CNN accelerator menunjukkan bahwa sistem mampu menghasilkan keluaran inferensi dalam 3802 siklus clock. Pada frekuensi kerja 100 MHz, nilai tersebut setara dengan latensi sebesar 38,02 $\mu$s.

    \item Hasil penggunaan resource menunjukkan bahwa CNN accelerator menggunakan sebagian besar DSP pada FPGA. Pada implementasi sistem penuh, penggunaan DSP mencapai 235 DSP48E1 dari total 240 DSP yang tersedia, sehingga DSP menjadi resource utama yang membatasi pengembangan model.

    \item Sistem streaming citra berbasis OV7670, DDR3 SDRAM, dan VGA berhasil dirancang dan diintegrasikan sebagai bagian dari sistem pendukung akuisisi dan tampilan citra pada FPGA.

    \item Hasil implementasi sistem penuh menunjukkan bahwa desain dapat melalui proses routing tanpa unrouted net dan memenuhi timing dengan nilai WNS positif. Estimasi daya total on-chip yang diperoleh adalah sebesar 1,099 W.
\end{enumerate}

\section{Saran}
Berdasarkan penelitian yang telah dilakukan, terdapat beberapa saran untuk pengembangan selanjutnya sebagai berikut.

\begin{enumerate}
    \item Optimasi penggunaan DSP perlu dilakukan karena implementasi CNN accelerator menggunakan resource DSP dalam jumlah sangat besar. Optimasi dapat dilakukan melalui penggunaan time multiplexing, pengurangan jumlah channel, atau perancangan ulang arsitektur CNN yang lebih ringan.

    \item Pengujian sistem dapat dikembangkan menggunakan data citra yang diperoleh langsung dari kamera OV7670 agar sistem klasifikasi lebih merepresentasikan kondisi penggunaan secara real-time.

    \item Pengukuran daya dapat dilakukan menggunakan alat ukur eksternal agar hasil konsumsi daya tidak hanya bergantung pada estimasi dari Vivado.

\end{enumerate}
